\documentclass[12pt]{article}
\usepackage[utf8]{inputenc}
\usepackage[russian]{babel}
\usepackage{amsmath, amssymb}
\usepackage{geometry}
\usepackage{parskip}

\usepackage{fancyhdr}
\usepackage[
  colorlinks=true,
  linkcolor=black,
  citecolor=black,
  urlcolor=blue,
  pdfborder={0 0 0}
]{hyperref}


\usepackage{fontspec}
\usepackage{setspace}

% Настройки страницы
\geometry{a4paper, margin=2cm}
\setlength{\parindent}{0pt}
\singlespacing
\pagestyle{empty}

% Шрифт (если xelatex)
\setmainfont{Liberation Serif} % или оставьте по умолчанию

\begin{document}

\begin{center}
\LARGE\textbf{M $\equiv$ F}

\bigskip

\large\textit{Ток как акт тождества между средой и полем}

\bigskip

> \large Ток — не движение зарядов. \\
> \large Ток — акт тождества между материальной системой и электрическим полем.

\bigskip

\[
\boxed{M \equiv F}
\]

\bigskip

Когда вы подключаете источник — \\
вы не "пускаете ток". \\
Вы нарушаете покой: \( M \not\equiv F \).  

Ответ системы — мгновенный. \\
Мы его называем \textbf{током}.

Это — не модель. \\
Это — онтология.

\bigskip\bigskip

\textbf{Следствия}

\medskip

\begin{tabular}{@{}l@{}}
• \textbf{Ом}: $ U $ — мера нарушения, $ R $ — способность восстановить. \\
• \textbf{Кирхгоф}: сумма = 0 — не баланс, а синхронный ответ. \\
• \textbf{Фарадей}: ЭДС — не наведение, а акт восстановления поля.
\end{tabular}

\medskip

Все они — проявления одного.


\bigskip

> Не кричу. \\
> Считаю. \\
> Утверждаю: это и есть обоснование.

\bigskip

Автор: Ивин Сергей Юрьевич \\
\texttt{git@ok-ru.ru}

\medskip

\href{https://github.com/OKRUgit/ivin-reflex-2026}{Исходный код и документы на }{\texttt{→ GitHub}}

\end{center}



\end{document}
